\documentclass[a4paper,11pt]{ltjsarticle}
\usepackage{luatexja}
\usepackage{luatexja-fontspec}
\usepackage{amsmath,amssymb,amsthm}
\usepackage{mathtools}
\usepackage{booktabs}
\usepackage{longtable}
\usepackage{array}
\usepackage{hyperref}
\usepackage{graphicx}
\usepackage{float}
\usepackage{geometry}
\geometry{margin=25mm}

\hypersetup{
  colorlinks=true,
  linkcolor=blue,
  urlcolor=blue,
  citecolor=blue
}

\theoremstyle{definition}
\newtheorem{theorem}{定理}[section]
\newtheorem{proposition}[theorem]{命題}
\newtheorem{lemma}[theorem]{補題}
\newtheorem{corollary}[theorem]{系}
\newtheorem{definition}[theorem]{定義}
\newtheorem{remark}[theorem]{注意}
\newtheorem{observation}[theorem]{観察}

\setlength{\parindent}{1em}
\setlength{\parskip}{0.5em}

\providecommand{\tightlist}{\setlength{\itemsep}{0pt}\setlength{\parskip}{0pt}}
\begin{document}

\section{背景空間と主観空間への次元追加についての考察}\label{ux80ccux666fux7a7aux9593ux3068ux4e3bux89b3ux7a7aux9593ux3078ux306eux6b21ux5143ux8ffdux52a0ux306bux3064ux3044ux3066ux306eux8003ux5bdf}

\begin{center}\rule{0.5\linewidth}{0.5pt}\end{center}

\textbf{著者}: 木原 範昭 (WF System Co., Ltd.)

\textbf{日付}: 2026年4月

\textbf{DOI：} \href{https://doi.org/10.5281/zenodo.19526913}{10.5281/zenodo.19526913}

\begin{center}\rule{0.5\linewidth}{0.5pt}\end{center}

\subsection{要旨}\label{ux8981ux65e8}

本稿では、先行論文 {[}1{]}, {[}2{]}, {[}3{]} で確立された中心投影フレームワーク、および論文 {[}4{]} で示された主観空間の入れ子構造の存在命題を前提として、背景空間および有限曲率を持つ主観空間への\textbf{新規次元の追加}について考察する。論文 {[}4{]} では入れ子構造の幾何学的存在可能性が示されたが、追加される次元が整合的に解釈可能である条件は別に検討を要する。本稿では、\textbf{追加される次元がどの曲率半径の支配下にあると解釈されるか} は観測者の解釈に依存する、という論文 {[}2{]} 定理 3.1 の延長線上にある観察から出発し、親空間の曲率半径 \(R\) の値(有限か、\(\infty\) か)によって、追加可能な次元の種類に構造的制約が生じることを示す。本稿で述べるのは、これらの制約の幾何学的・論理的な記述のみであり、その物理的実在性については何も主張しない。

\begin{center}\rule{0.5\linewidth}{0.5pt}\end{center}

\subsection{1. 導入}\label{ux5c0eux5165}

先行論文 {[}1{]} で導入された中心投影フレームワークにおいて、\((n+1)\) 次元ユークリッド背景空間 \(\mathbb{R}^{n+1}\) の中に半径 \(R\) の超球面 \(S^n(R)\) および接超平面 \(\Pi_R\) が定義される。主観空間 \(\Pi_R\) は曲率半径 \(R\) に支配される \(n\) 次元の空間であり、Einstein 方程式 \(G_{\mu\nu} + \Lambda g_{\mu\nu} = 0\), \(\Lambda = 3/R^2\) を満たす。

論文 {[}2{]} 定理 3.1 は、主観空間の軸の間に幾何学的な絶対的対等性があり、どの軸を特定の役割(例えば時間に対応する軸)と解釈するかは観測者の選択に委ねられることを示した。

論文 {[}4{]} §4 では、任意の主観空間の内部に新たな主観空間を立ち上げる自由度が、既存のフレームワークから直接導出されることが示された(命題 4.1: 入れ子構造の存在)。

しかし、論文 {[}4{]} の命題 4.1 は「立ち上げが可能である」という存在命題であり、立ち上げられた構造が\textbf{整合的に解釈可能である条件}は別途検討を要する。本稿の目的は、この解釈の整合性条件を、親空間の曲率半径 \(R\) の値に応じて検討することにある。

\begin{center}\rule{0.5\linewidth}{0.5pt}\end{center}

\subsection{2. 次元追加の解釈依存性}\label{ux6b21ux5143ux8ffdux52a0ux306eux89e3ux91c8ux4f9dux5b58ux6027}

\subsubsection{2.1 解釈依存性の一般的性質}\label{ux89e3ux91c8ux4f9dux5b58ux6027ux306eux4e00ux822cux7684ux6027ux8cea}

論文 {[}2{]} 定理 3.1 において、主観空間の軸の間の対等性は、どの軸が「時間に対応する」かといった個別的解釈が観測者の選択に委ねられることを含意する。本稿ではこの解釈依存性を、次元追加の場面に拡張する。

主観空間または背景空間に新規の次元を追加するとき、その新規次元が\textbf{どの曲率半径の支配下にあると解釈されるか} は、観測者の解釈に依存する。具体的には、新規次元が既存の曲率半径 \(R\) の支配下にあると解釈する立場と、新たな独立した曲率半径 \(R'\) の支配下にあると解釈する立場の、いずれも許容されうる。

幾何学的な追加操作それ自体は共通であっても、その解釈が異なれば、新規次元が意味する構造は異なる。本稿で記述する「制約」は、幾何学的な禁止ではなく、解釈の整合性から生じる論理的制約である。

\subsubsection{2.2 二種類の次元追加}\label{ux4e8cux7a2eux985eux306eux6b21ux5143ux8ffdux52a0}

以上の議論から、次元追加には以下の二種類が区別される。

\textbf{種類 I}(支配下の追加)　新規次元が、既存の曲率半径 \(R\) の支配下にあると解釈される場合。この場合、新規次元は既存の主観空間の拡張として理解され、新たな曲率構造を導入しない。

\textbf{種類 II}(新曲率半径の導入)　新規次元が、新たな独立した曲率半径 \(R'\) の支配下にあると解釈される場合。この場合、新規次元は別の主観空間 \(\Pi_{R'}\) の立ち上げに対応する。

以下、親空間の曲率半径 \(R\) の値によって、これら二種類の追加が許容されるか否かを検討する。

\begin{center}\rule{0.5\linewidth}{0.5pt}\end{center}

\subsection{3. 平坦な背景空間への次元追加}\label{ux5e73ux5766ux306aux80ccux666fux7a7aux9593ux3078ux306eux6b21ux5143ux8ffdux52a0}

\subsubsection{3.1 状況の設定}\label{ux72b6ux6cc1ux306eux8a2dux5b9a}

背景空間 \(\mathbb{R}^{n+1}\) は曲率を持たない。これはフレームワーク全体の出発点であり、\(R = \infty\) に対応する状態とみなすことができる。

この背景空間に対して、新規次元を追加する場合を検討する。

\subsubsection{3.2 許容される追加}\label{ux8a31ux5bb9ux3055ux308cux308bux8ffdux52a0}

\textbf{観察 3.1}　平坦な背景空間 \(\mathbb{R}^{n+1}\) に対して、新規次元が整合的に解釈可能である条件は、以下のいずれかを満たすことである。

\textbf{(a)} 新規次元が \(R_1 = \infty\) に支配されていると解釈される。すなわち、追加後も背景空間の平坦性は維持される。

\textbf{(b)} 新規次元が、既存または新規の有限曲率半径に支配されていると解釈され、かつその解釈が背景空間の住民の視点から整合的である。

条件 (a) は、背景空間の次元を増やす操作に相当し、平坦性を保ちつつ \((n+1) \to (n+2)\) へと次元を拡張することに等しい。この場合、幾何学的にも解釈的にも整合性は自明に成立する。

条件 (b) は、追加される次元が背景空間の住民から見て「ある有限曲率半径の一部」として解釈可能である場合にのみ許容される。この解釈が成立しない場合、すなわち新規次元が背景空間の平坦性と独立した曲率構造を導入すると解釈される場合、論文 {[}2{]} 定理 3.1 の軸の対等性と組み合わせると、背景空間全体に当該曲率構造が遍在することになり、出発点である背景空間の平坦性と矛盾する。

\subsubsection{3.3 要点}\label{ux8981ux70b9}

平坦な背景空間に対して、種類 I の追加は常に許容される(条件 (a))。種類 II の追加は、その新規曲率半径に支配される構造が背景空間の住民の視点から整合的に解釈できる場合に限って許容される(条件 (b))。

\begin{center}\rule{0.5\linewidth}{0.5pt}\end{center}

\subsection{4. 有限曲率の主観空間への次元追加}\label{ux6709ux9650ux66f2ux7387ux306eux4e3bux89b3ux7a7aux9593ux3078ux306eux6b21ux5143ux8ffdux52a0}

\subsubsection{4.1 状況の設定}\label{ux72b6ux6cc1ux306eux8a2dux5b9a-1}

主観空間 \(\Pi_R\) は、有限の曲率半径 \(R < \infty\) によって支配される \(n\) 次元空間である。この主観空間に対して、新規次元を追加する場合を検討する。

\subsubsection{4.2 種類 I の追加: 既存の曲率半径の支配下}\label{ux7a2eux985e-i-ux306eux8ffdux52a0-ux65e2ux5b58ux306eux66f2ux7387ux534aux5f84ux306eux652fux914dux4e0b}

\textbf{観察 4.1}　主観空間 \(\Pi_R\) に対して、新規次元が既存の曲率半径 \(R\) の支配下にあると解釈される場合、この種類 I の追加は\textbf{任意に実行可能}である。

追加される次元は、既存の主観空間 \(\Pi_R\) と同じ \(R\) を共有するので、新たな曲率構造を導入しない。論文 {[}2{]} 定理 3.1 の軸の対等性により、既存の軸と追加された軸の区別は観測者の解釈の問題であり、幾何学的には対等である。したがって、追加可能な次元の数に原理的な制約はない。

\subsubsection{4.3 種類 II の追加: 新曲率半径の導入}\label{ux7a2eux985e-ii-ux306eux8ffdux52a0-ux65b0ux66f2ux7387ux534aux5f84ux306eux5c0eux5165}

\textbf{観察 4.2}　主観空間 \(\Pi_R\) に対して、新規次元が新たな独立した曲率半径 \(R' \neq R\) の支配下にあると解釈される場合、論文 {[}4{]} 命題 4.1 によりその立ち上げは幾何学的に可能である。しかし、この新規次元\textbf{単独}の追加では、立ち上げられる主観空間 \(\Pi_{R'}\) の構造は退化する。

具体的には、新規曲率半径 \(R'\) の軸それ自身は \(\Pi_{R'}\) の外部方向に対応し、\(\Pi_{R'}\) の住民からは認知されない。したがって \(\Pi_{R'}\) の住民から見える次元は、\(R'\) の支配下にあると解釈される追加的な次元によって構成される必要がある。

この追加的な次元が同時に追加されない場合、\(\Pi_{R'}\) の住民から見た \(\Pi_{R'}\) の次元数は \(0\) となる。すなわち、\(\Pi_{R'}\) は\textbf{ゼロ次元の主観空間}となる。

\textbf{観察 4.3}　したがって、種類 II の追加が非退化な主観空間 \(\Pi_{R'}\) を構成するためには、新規曲率半径 \(R'\) に対応する軸と、\(R'\) の支配下にあると解釈される追加的な次元を、\textbf{同時に追加}する必要がある。これらの次元は種類 I の追加と同じ形式で、\(R'\) の支配下に任意の数だけ追加可能である。

\subsubsection{4.4 要点}\label{ux8981ux70b9-1}

有限曲率の主観空間に対して、種類 I の追加は任意の数だけ実行可能である。種類 II の追加は、新曲率半径 \(R'\) に支配される追加的な次元を同時に追加しない限り、立ち上げられる新主観空間はゼロ次元となる。

\begin{center}\rule{0.5\linewidth}{0.5pt}\end{center}

\subsection{5. 考察}\label{ux8003ux5bdf}

本稿では、先行論文 {[}1{]}, {[}2{]}, {[}3{]}, {[}4{]} のフレームワークに基づき、次元追加の解釈的整合性条件を記述した。

本稿の主張は以下に限定される。

\begin{enumerate}
\def\labelenumi{\arabic{enumi}.}
\tightlist
\item
  次元追加の種類は、解釈依存的に二種類(種類 I: 支配下の追加、種類 II: 新曲率半径の導入)に区別される。
\item
  平坦な背景空間に対しては、種類 I は常に許容され、種類 II は解釈が背景空間の住民の視点から整合的である場合に限って許容される。
\item
  有限曲率の主観空間に対しては、種類 I は任意に許容され、種類 II は新曲率半径の支配下にある追加的な次元を同時に追加しない限り、立ち上げられる新主観空間はゼロ次元となる。
\end{enumerate}

これらの制約は幾何学的な絶対的禁止ではなく、観測者の解釈の整合性から生じる論理的制約である。この意味で、本稿の制約は論文 {[}2{]} 定理 3.1 の軸の対等性(時間軸に対応する軸の選択が観測者の解釈に依存する)の自然な延長線上にある。

本稿はこれらの制約の物理的実在性については何も主張しない。また、ゼロ次元となる新主観空間の物理的意味付け、および種類 II の追加における「支配下の次元を同時に追加する」という条件の物理的含意についても、本稿では踏み込まない。これらの問いは幾何学の外側の問題であり、読者の判断に委ねられる。

\begin{center}\rule{0.5\linewidth}{0.5pt}\end{center}

\subsection{6. 結論}\label{ux7d50ux8ad6}

中心投影フレームワークにおける次元追加は、解釈依存的に二種類に区別される。その各々が整合的に解釈可能である条件は、親空間の曲率半径 \(R\) の値に依存する。平坦な背景空間 (\(R = \infty\)) では種類 II の追加に背景空間の住民からの整合性という制約が課される。有限曲率の主観空間 (\(R < \infty\)) では種類 I は任意に許容され、種類 II は支配下の次元の同時追加を要する。物理的解釈は本稿の対象外であり、読者に委ねる。

\begin{center}\rule{0.5\linewidth}{0.5pt}\end{center}

\subsection{参考文献}\label{ux53c2ux8003ux6587ux732e}

{[}1{]} 木原 範昭,「中心投影による幾何学的定式化」, Zenodo, DOI: 10.5281/zenodo.19427780 (2025).

{[}2{]} 木原 範昭,「中心投影フレームワークにおける対称性」, Zenodo, DOI: 10.5281/zenodo.19434932 (2025).

{[}3{]} 木原 範昭,「複数主観空間における観測の相対性」, Zenodo, DOI: 10.5281/zenodo.19435162 (2025).

{[}4{]} 木原 範昭,「中心投影における曲率半径の極限についての考察 --- \(R \to \infty\) と \(R \to 0\) の場合」, Zenodo, DOI: 10.5281/zenodo.19526549 (2026).

\begin{center}\rule{0.5\linewidth}{0.5pt}\end{center}

\end{document}
